\documentclass[11pt]{article}
\usepackage[utf8]{inputenc}
\usepackage[T1]{fontenc}
\usepackage[spanish]{babel}
\usepackage{amsmath,amssymb}
\usepackage{geometry}
\usepackage{booktabs}
\usepackage{array}
\usepackage{hyperref}
\usepackage{xcolor}
\usepackage{graphicx}
\usepackage{float}
\usepackage{tikz}
\usetikzlibrary{arrows.meta,positioning,shapes.geometric,calc}

\geometry{letterpaper, margin=1in}
\hypersetup{colorlinks=true, linkcolor=black, urlcolor=black, hypertexnames=false}
\newcolumntype{L}[1]{>{\raggedright\arraybackslash}p{#1}}

\definecolor{gOne}{gray}{0.93}
\definecolor{gTwo}{gray}{0.86}
\definecolor{gThree}{gray}{0.78}
\definecolor{gFour}{gray}{0.70}

\tikzset{
  box/.style={rectangle, rounded corners=2mm, draw=black!70, very thick, align=center, minimum height=9mm, font=\small},
  sbox/.style={rectangle, rounded corners=2mm, draw=black!70, thick, align=center, minimum height=8mm, font=\scriptsize},
  arrow/.style={-{Latex[length=2.2mm]}, thick, draw=black!70},
  softarrow/.style={-{Latex[length=2.2mm]}, thick, dashed, draw=black!55}
}

\title{El tatuaje como archivo narrativo:\\ lectura situada de fuentes digitales mediante un sistema local}
\author{}
\date{}

\begin{document}
\maketitle

\begin{abstract}
El tatuaje puede leerse como archivo corporal, marca estética, memoria afectiva,
oficio visual y objeto de disputa social. El análisis de fuentes digitales
enfrenta, sin embargo, un problema metodológico frecuente: reunir noticias,
foros, reportes y artículos académicos puede producir un corpus grande, pero
también una voz artificialmente homogénea. Este trabajo presenta un
procedimiento local para construir un corpus narrativo trazable y analizarlo sin
enviar textos a servicios externos. El ejemplo conductor es el tópico
\textit{tatuaje}, entendido no sólo como palabra clave, sino como campo de
sentidos: cuerpo, identidad, oficio, estética, salud, regulación, memoria y
estigma. El sistema separa rubros de búsqueda, años y capas de fuente; depura
los textos; registra vacíos y fallas; permite organizar indicios narrativos,
construir mapas de relaciones y proponer síntesis revisables para lectura
cercana. No se
busca automatizar la interpretación, sino producir una mediación metodológica
auditable para historiadores, artistas, diseñadores y estudios culturales.
\end{abstract}

\section{Problema cultural y metodológico}

Una narrativa social no aparece en una sola clase de documento. Puede emerger en
conversaciones cotidianas, circular por la prensa, ser traducida a lenguaje
institucional y después estabilizarse en estudios académicos. En el caso del
tatuaje, una misma palabra puede remitir a prácticas artísticas, marcas
comerciales, procedimientos médicos, identidades juveniles, estigmas laborales o
regulaciones sanitarias. Si todo se busca junto, el corpus puede crecer, pero
perder sentido.

El objetivo del procedimiento es construir un corpus situado que conserve tres cosas:
procedencia, diferencia entre voces y decisiones de limpieza. La pregunta no es
``qué dice Internet sobre el tatuaje'', sino cómo distintas capas discursivas
producen, desplazan o disputan sentidos alrededor del tatuaje en un periodo y
región definidos.

\section{Fases del procedimiento}

\begin{table}[H]
\centering
\small
\begin{tabular}{L{0.16\textwidth}L{0.36\textwidth}L{0.38\textwidth}}
\toprule
\textbf{Fase} & \textbf{Qué se hace} & \textbf{Por qué se hace} \\
\midrule
Observar el campo & Se reconoce el tatuaje como inscripción corporal, práctica visual y relato social. &
Sitúa el problema antes de convertirlo en búsqueda digital. \\
\addlinespace
Delimitar & Se define región, periodo y rubros de sentido: identidad, oficio, salud, memoria, diseño. &
Evita mezclar sentidos distintos bajo una misma palabra. \\
\addlinespace
Reunir voces & Se consulta por capas: prensa, foros, artículos y reportes. &
Permite distinguir conversación pública, prensa, investigación y documentos institucionales. \\
\addlinespace
Cuidar el corpus & Se eliminan duplicados, menús, publicidad, ruido y homónimos. &
Impide que el ruido de páginas web se vuelva resultado interpretativo. \\
\addlinespace
Interpretar relaciones & Se observan expresiones, actores, momentos narrativos, fuentes y tonalidad léxica. &
Produce indicios y mapas de sentido para lectura crítica. \\
\addlinespace
Síntesis crítica & Se proponen nodos relevantes y se vuelve a ejemplos textuales. &
Prepara una lectura cercana sin presentar el algoritmo como interpretación final. \\
\bottomrule
\end{tabular}
\caption{Fases metodológicas para una lectura situada del tatuaje en fuentes digitales.}
\label{tab:fases-es}
\end{table}

\begin{figure}[H]
\centering
\resizebox{0.96\textwidth}{!}{%
\begin{tikzpicture}[node distance=7mm]
  \node[box, fill=gOne, minimum width=31mm] (a) {1. Situar\\archivo corporal};
  \node[box, fill=gTwo, minimum width=31mm, right=of a] (b) {2. Reunir voces\\capas discursivas};
  \node[box, fill=gThree, minimum width=31mm, right=of b] (c) {3. Cuidar corpus\\limpiar y auditar};
  \node[box, fill=gFour, minimum width=31mm, right=of c] (d) {4. Leer relaciones\\síntesis crítica};
  \draw[arrow] (a) -- (b);
  \draw[arrow] (b) -- (c);
  \draw[arrow] (c) -- (d);
  \draw[softarrow] (d.south) .. controls +(0,-10mm) and +(0,-10mm) .. node[below,font=\scriptsize] {revisión humana: ajustar rubros, exclusiones y lectura} (a.south);
\end{tikzpicture}%
}
\caption{Itinerario interpretativo: el sistema ordena la lectura, no sustituye al investigador.}
\label{fig:fases-es}
\end{figure}

\section{Ejemplo conductor: tatuaje}

Para un público de historia, arte, diseño o estudios culturales, el tatuaje no
debe tratarse como una palabra aislada. Conviene dividirlo en rubros:

\begin{table}[H]
\centering
\small
\begin{tabular}{L{0.24\textwidth}L{0.62\textwidth}}
\toprule
\textbf{Rubro} & \textbf{Variantes de búsqueda} \\
\midrule
Núcleo & tatuaje, tatuajes, tattoo, tattoos, arte corporal \\
Oficio e industria & tatuador, tatuadora, estudio de tatuajes, artista del tatuaje \\
Identidad y memoria & significado de tatuaje, tatuaje e identidad, tatuaje y memoria, tatuaje religioso \\
Sociedad y trabajo & tatuaje y juventud, tatuaje y empleo, tatuaje y discriminación, tatuaje y género \\
Salud y regulación & tatuaje y salud, tintas para tatuaje, infección, regulación sanitaria \\
Estética y diseño & diseño de tatuaje, tatuaje tradicional, tatuaje mexicano, body art \\
\bottomrule
\end{tabular}
\caption{Rubros iniciales para el tópico tatuaje. Los rubros son editables y deben validarse por lectura.}
\label{tab:tatuaje-rubros-es}
\end{table}

También se declaran exclusiones. En este caso son necesarias porque
\textit{Tatuaje} puede aparecer como marca de cigarros o como procedimiento
médico especializado. Excluir \textit{cigar}, \textit{tobacco}, \textit{robusto}
o \textit{colonoscopic tattooing} no censura el corpus: protege la pregunta de
investigación.

\begin{figure}[H]
\centering
\resizebox{\textwidth}{!}{%
\begin{tikzpicture}[node distance=9mm]
  \node[box, fill=gOne, minimum width=32mm] (q) {Pregunta\\práctica cultural};
  \node[box, fill=gOne, minimum width=31mm, right=of q] (r) {Rubros\\sentidos};
  \node[box, fill=gTwo, minimum width=32mm, right=of r] (y) {Años y región\\México, 2020--2026};
  \node[box, fill=gTwo, minimum width=34mm, below=of q] (v) {Capas de voz\\prensa, foros, academia};
  \node[box, fill=gThree, minimum width=34mm, right=of v] (c) {Cuidado\\limpiar y auditar};
  \node[box, fill=gFour, minimum width=34mm, right=of c] (m) {Lectura\\relaciones};
  \draw[arrow] (q) -- (r);
  \draw[arrow] (r) -- (y);
  \draw[arrow] (q) -- (v);
  \draw[arrow] (r) -- (v);
  \draw[arrow] (y) -- (v);
  \draw[arrow] (v) -- (c);
  \draw[arrow] (c) -- (m);
\end{tikzpicture}%
}
\caption{Diseño de corpus para el caso tatuaje: observe la separación entre pregunta cultural, rubros de sentido y capas de voz antes de reunir documentos.}
\label{fig:corpus-tatuaje-es}
\end{figure}

\section{De documentos a huellas narrativas}

El sistema guarda una base por cada combinación de rubro, año y capa de fuente.
Si una búsqueda no produce resultados, el vacío se registra. Esto importa: una
ausencia también informa sobre el alcance del corpus y sobre los límites de la
recolección.

La recolección actual usa una estrategia secuencial adaptativa: primero se
consulta el término base del rubro y sólo se expande a sinónimos si hace falta.
Esto evita saturar los índices y permite saber qué palabra activó cada hallazgo.
Para el caso del tatuaje, no es lo mismo recuperar documentos por
\textit{tatuaje}, \textit{tatuajes}, \textit{tattoo} o \textit{arte corporal};
cada término arrastra comunidades, idiomas y sesgos distintos.

\begin{table}[H]
\centering
\small
\begin{tabular}{L{0.18\textwidth}L{0.25\textwidth}L{0.42\textwidth}}
\toprule
\textbf{Paso} & \textbf{Ejemplo} & \textbf{Lectura metodológica} \\
\midrule
Entrada & ``tatuaje y salud'', México, 2021, noticias & Se acota un sentido del tópico en una capa pública. \\
Depuración & Se retiran menús, anuncios, duplicados y homónimos & El corpus no confunde publicidad o ruido con narrativa social. \\
Descripción & Aparecen términos como tinta, piel, regulación, riesgo & Se producen indicios, no conclusiones. \\
Relación & Texto--actor--idea--fuente--año & Se observa cómo se conectan voces y conceptos. \\
Síntesis & Nodo relevante: regulación sanitaria / estudio de tatuajes & Se elige un punto de entrada para lectura cercana. \\
\bottomrule
\end{tabular}
\caption{Mini-demostración del procedimiento con el rubro salud y regulación.}
\label{tab:mini-demo-es}
\end{table}

\section{Derechos de fuente y grados de evidencia}

La extracción debe respetar los límites de cada fuente. El sistema no intenta
entrar a espacios privados, no automatiza sesiones y no evade restricciones. Si
una página pública permite extracción según \texttt{robots.txt}, se intenta
obtener texto completo. Si no lo permite, si el sitio es parcial o si hay
paywall, se conserva sólo una señal trazable: título, resumen, medio, fecha,
URL y tipo de fuente. Esa señal puede orientar el mapa del corpus, pero no debe
contarse como cuerpo textual equivalente.

\begin{figure}[H]
\centering
\resizebox{0.98\textwidth}{!}{%
\begin{tikzpicture}[node distance=9mm]
  \node[box, fill=gOne, minimum width=34mm] (i) {Índice público\\RSS, GDELT, OpenAlex};
  \node[box, fill=gTwo, minimum width=34mm, right=of i] (u) {URL o metadato\\fuente, año, término};
  \node[box, fill=gThree, minimum width=34mm, right=of u] (p) {Permiso\\robots.txt / acceso};
  \node[box, fill=gFour, minimum width=34mm, below left=of p] (ok) {Texto completo\\estado ok};
  \node[box, fill=gFour, minimum width=34mm, below right=of p] (parcial) {Señal parcial\\estado ok\_partial};
  \node[box, fill=gTwo, minimum width=38mm, below=of ok] (local) {Análisis local\\sin redistribuir textos};
  \draw[arrow] (i) -- (u);
  \draw[arrow] (u) -- (p);
  \draw[arrow] (p) -- node[left,font=\scriptsize] {permitido} (ok);
  \draw[arrow] (p) -- node[right,font=\scriptsize] {parcial/no permitido} (parcial);
  \draw[arrow] (ok) -- (local);
  \draw[arrow] (parcial) -- (local);
\end{tikzpicture}%
}
\caption{Grados de evidencia: texto completo y señal parcial se guardan separados para no exagerar el alcance del corpus.}
\label{fig:derechos-evidencia-es}
\end{figure}

Esta diferencia es central para estudiar lo que las personas dicen en red. Los
foros, Reddit público, blogs y comunidades abiertas aportan señales de
conversación pública, pero no representan toda la conversación social. Deben
leerse como huellas indexables, con sesgos de plataforma, idioma, moderación,
visibilidad y acceso.

\subsection*{Microcaso de lectura}

Una corrida publicable no debe quedarse en el conteo de documentos. Debe mostrar
al menos una cadena mínima: fuente, fragmento limpiado, etiqueta asignada,
relación construida y lectura humana. En el caso del tatuaje, una nota de prensa
sobre tintas o regulación sanitaria puede conectar autoridades de salud,
estudios de tatuaje y vocabulario de riesgo; una conversación en foro puede
conectar clientes, dolor, identidad, precio o confianza; un artículo académico
puede estabilizar vocabulario médico sobre piel, infección o cicatrización.

\begin{table}[H]
\centering
\small
\begin{tabular}{L{0.17\textwidth}L{0.24\textwidth}L{0.23\textwidth}L{0.24\textwidth}}
\toprule
\textbf{Capa} & \textbf{Indicio posible} & \textbf{Nodo/relación} & \textbf{Lectura humana} \\
\midrule
Prensa & ``regulación de tintas'' & autoridad sanitaria--tinta--riesgo & El tatuaje aparece como asunto público y sanitario. \\
Foro & ``recomendación de estudio'' & cliente--tatuador--confianza & El oficio se evalúa por experiencia compartida. \\
Artículo & ``infección cutánea'' & piel--procedimiento--cuidado & El cuerpo se vuelve objeto de saber médico. \\
\bottomrule
\end{tabular}
\caption{Ejemplo de cómo un documento se transforma en entrada de lectura, no en conclusión automática.}
\label{tab:microcaso-es}
\end{table}

\section{Mapas de relaciones}

La red narrativa se usa como mapa. Cada texto conecta actores, ideas, fuente,
año, lugar y momentos del relato. El mapa no prueba por sí mismo una hipótesis;
ayuda a decidir dónde leer con más atención.

Los actores, ideas y momentos del relato no deben asumirse como entidades
transparentes. Pueden provenir de reglas locales, diccionarios editables,
frecuencias de n-gramas, patrones lingüísticos y validación humana. En un tema
como tatuaje, los actores esperables incluyen tatuadores, clientes, autoridades
sanitarias, médicos, empleadores, medios y colectivos juveniles; las ideas pueden
agrupar identidad, memoria, oficio, riesgo, estética, discriminación o cuidado.
Cada asignación debe poder corregirse desde la lectura.

\begin{figure}[H]
\centering
\begin{tikzpicture}[node distance=17mm]
  \node[circle, draw=black!70, very thick, fill=gTwo, minimum size=15mm, align=center] (texto) {Texto};
  \node[circle, draw=black!70, thick, fill=gOne, above left=of texto, minimum size=14mm] (actor) {Actor};
  \node[circle, draw=black!70, thick, fill=gOne, above right=of texto, minimum size=17mm, align=center] (momento) {Momento\\del relato};
  \node[circle, draw=black!70, thick, fill=gThree, right=of texto, minimum size=14mm] (idea) {Idea};
  \node[circle, draw=black!70, thick, fill=gOne, below left=of texto, minimum size=14mm] (fuente) {Fuente};
  \node[circle, draw=black!70, thick, fill=gOne, below=of texto, minimum size=12mm] (ano) {Año};
  \node[circle, draw=black!70, thick, fill=gOne, below right=of texto, minimum size=14mm] (lugar) {Lugar};
  \draw[arrow] (texto) -- node[left,font=\scriptsize] {menciona} (actor);
  \draw[arrow] (texto) -- node[right,font=\scriptsize] {contiene} (momento);
  \draw[arrow] (texto) -- node[above,font=\scriptsize] {usa} (idea);
  \draw[arrow] (fuente) -- node[left,font=\scriptsize] {publica} (texto);
  \draw[arrow] (ano) -- node[right,font=\scriptsize] {fecha} (texto);
  \draw[arrow] (lugar) -- node[right,font=\scriptsize] {sitúa} (texto);
  \draw[arrow] (actor) -- node[above,font=\scriptsize] {participa en} (momento);
  \draw[arrow] (momento) -- node[right,font=\scriptsize] {asocia} (idea);
\end{tikzpicture}
\caption{Red narrativa mínima: el lector ve qué tipo de relación está siendo registrada.}
\label{fig:red-es}
\end{figure}

\section{Tonalidad léxica exploratoria}

La tonalidad léxica no se interpreta como emoción colectiva ni como estado
psicológico de quienes escriben. Se usa como indicador exploratorio de
vocabulario valorativo. Por ejemplo, puede mostrar si las noticias hablan del
tatuaje con términos de riesgo, si los foros usan términos de identidad o si los
artículos científicos concentran vocabulario de salud y regulación.

El indicador básico compara palabras positivas y negativas, con reglas simples
de negación e intensificación. El resultado se resume por año y por capa de
fuente. En una gráfica radial, cada eje representa una capa: prensa, foros,
artículos o reportes. La interpretación siempre debe volver a ejemplos
textuales.

\section{Síntesis revisable}

Cuando el mapa tiene muchos elementos, el sistema propone nodos relevantes:
actores, ideas, fuentes o momentos del relato que concentran relaciones. Esta
selección no reemplaza la lectura. Sirve para preguntar qué entradas conviene
leer primero y qué relaciones se perderían si se retiraran ciertos nodos.

\begin{figure}[H]
\centering
\resizebox{0.96\textwidth}{!}{%
\begin{tikzpicture}[node distance=8mm]
  \node[box, fill=gOne, minimum width=34mm] (corpus) {Corpus depurado\\textos y procedencia};
  \node[box, fill=gTwo, minimum width=34mm, right=of corpus] (mapa) {Mapa de relaciones\\actores, ideas, fuentes};
  \node[box, fill=gThree, minimum width=34mm, right=of mapa] (nodos) {Nodos relevantes\\posibles entradas};
  \node[box, fill=gFour, minimum width=36mm, right=of nodos] (lectura) {Lectura cercana\\volver al texto};
  \draw[arrow] (corpus) -- (mapa);
  \draw[arrow] (mapa) -- (nodos);
  \draw[arrow] (nodos) -- (lectura);
  \draw[softarrow] (lectura.south) .. controls +(0,-11mm) and +(0,-11mm) .. node[below,font=\scriptsize] {ajustar exclusiones, rubros o interpretación} (corpus.south);
\end{tikzpicture}%
}
\caption{La síntesis es revisable: el algoritmo sugiere entradas, la interpretación vuelve al corpus.}
\label{fig:sintesis-es}
\end{figure}

\section{Salida reproducible}

La aplicación local exporta un archivo estructurado de evidencia que conserva
registros filtrados, procedencia, limpieza, clasificación de fuente, tonalidad
léxica, eventos narrativos, actores, grupos de ideas y mapas de relaciones. Para
una publicación robusta falta consolidar además un manifiesto de corrida con
fecha, parámetros, versión del código, cuotas, semillas y huella del corpus
final.

\section{Límites y sesgos}

El corpus no representa por sí mismo a una sociedad. Representa una colección
situada de huellas disponibles: lo que los índices encuentran, lo que las
plataformas permiten consultar, lo que los medios publican y lo que ciertos
grupos escriben en línea. Puede haber sesgo de idioma, clase social,
alfabetización digital, centralidad urbana, disponibilidad de fuentes,
moderación de plataformas, agenda periodística y sobrerrepresentación sanitaria
o institucional. Las exclusiones también son decisiones interpretativas: retirar
\textit{colonoscopic tattooing} es razonable si se estudia tatuaje como práctica
cultural, pero podría ser incorrecto si el rubro buscado es medicina o salud.
Por eso cada exclusión debe registrarse y justificarse.

\appendix
\section{Nota técnica mínima}

La selección de nodos relevantes puede formalizarse como una cobertura sobre el
mapa de relaciones. Se buscan conjuntos pequeños de nodos que mantengan alta
relevancia y bajo daño estructural. En términos operativos, se comparan tres
criterios: tamaño del conjunto, peso narrativo de los nodos y relaciones que se
perderían si esos nodos se retiraran. Los detalles algorítmicos quedan como
apéndice porque la pregunta principal es interpretativa: ¿qué conjunto de
actores, ideas o fuentes ofrece una entrada defendible para leer el corpus?

\end{document}
